\section{Setting and error anatomy of regression GATE estimators}\label{sec:theory}
\textbf{Setting.} A fixed, connected, undirected graph with adjacency $A$ has $n$ nodes with degrees $d_i$. A design draws $z\in\{0,1\}^n$. Unit $i$ has expected outcome $\mu_i(z)$ and we observe $y_i=\mu_i(z)+\varepsilon_i$ with independent noise, $\mathrm{Var}(\varepsilon_i)=\sigma^2$. The estimand is
$\tau=\frac1n\sum_i[\mu_i(\mathbf 1)-\mu_i(\mathbf 0)]$ (GATE). The population (graph, $\mu_i(\cdot)$) is fixed; randomness is in $z$ and $\varepsilon$.

\textbf{Regression estimator.} Let $\phi_i(z)\in\mathbb R^q$ be features built from the assignment and the (observed) graph, e.g.\ $\phi_i=(z_i,x_i)$ with $x_i=\frac1{d_i}\sum_jA_{ij}z_j$ the treated neighbour fraction. Fit $y$ on $(1,\phi_i(z))$ by least squares to get $\hat\beta$, and set $\hat\tau=\frac1n\sum_i[\phi_i(\mathbf 1)-\phi_i(\mathbf 0)]^\top\hat\beta_{\phi}$. With $X=[\mathbf 1,\Phi(z)]$ and $\tilde c=(0,\frac1n\sum_i(\phi_i(\mathbf 1)-\phi_i(\mathbf 0)))$:
\begin{equation}
\hat\tau=\tilde c^\top(X^\top X)^{-1}X^\top y=\textstyle\sum_iw_iy_i,\qquad w=X(X^\top X)^{-1}\tilde c.\label{eq:w}
\end{equation}
\begin{proposition}[Exact error decomposition]\label{prop:decomp}
$w$ depends only on $z$ and the graph, $\sum_iw_i=0$ and $X^\top w=\tilde c$. For \emph{any} $\beta$, writing $\mu_i(z')=\phi_i(z')^\top\beta_\phi+\beta_0+r_i(z')$,
\[
\hat\tau-\tau=\underbrace{\textstyle\sum_iw_i\varepsilon_i}_{\text{noise}}+\underbrace{\textstyle\sum_iw_ir_i(z)-\tfrac1n\textstyle\sum_i\big[r_i(\mathbf 1)-r_i(\mathbf 0)\big]}_{\text{misspecification / extrapolation}} .
\]
Given $z$, the noise term has mean $0$ and variance $\sigma^2\|w\|_2^2=\sigma^2\tilde c^\top(X^\top X)^{-1}\tilde c$ (exactly, for any $n$, no asymptotics). The second term vanishes if $\mu$ is in the model, and otherwise satisfies $|\cdot|\le\|w\|_1\|r(z)\|_\infty+|\frac1n\sum_i(r_i(\mathbf 1)-r_i(\mathbf 0))|$ and $|\sum_iw_ir_i|\le\|w\|_2\|r(z)\|_2$.
\end{proposition}
\emph{Proof.} Since $X^\top w=\tilde c$, $\sum_iw_i\mu_i(z)=\tilde c^\top\beta+\sum_iw_ir_i(z)$, while $\tau=\tilde c^\top\beta+\frac1n\sum_i(r_i(\mathbf 1)-r_i(\mathbf 0))$. Subtract and add $\sum_iw_i\varepsilon_i$. \hfill$\square$

So the realised assignment fixes two numbers that can be computed \emph{before any outcome is seen}: $\|w\|_2$ (sensitivity to noise) and $\|w\|_1$ or $\|w\|_2$ (sensitivity to misspecification). We use them as design diagnostics and test them in \S\ref{sec:res}. In the experiments the second term is measured directly by refitting on the noiseless outcomes $\mu(z)$: this gives the \emph{assignment-induced} part $S=\sum_iw_i\mu_i(z)-\tau$ (random only through $z$), whose mean is the bias and whose variance is the design variance. The total error is $S+N$ with $N=\sum_iw_i\varepsilon_i$.

\begin{proposition}[Bernoulli design, regular graph; large $n$]\label{prop:bern}
Let $A$ be $d$-regular with bounded $d$, $z_i\sim\mathrm{Bern}(p)$ i.i.d., $\phi=(z,x)$. Then $n\|w\|_2^2\to\frac{d+1}{p(1-p)}$ as $n\to\infty$. Hence $\mathrm{RMSE}_{\rm noise}\approx\sigma\sqrt{(d+1)/(np(1-p))}$: error grows like $\sqrt d$ although each unit has more neighbours. If $\mu_i(z)=a_i+\tau_dz_i+\gamma g(x_i)$ for a smooth $g$, and the $a_i$ are exogenous baseline terms, the population least-squares slope on $x$ is $\beta_x=\mathrm{Cov}(g(x),x)/\mathrm{Var}(x)=g'(p)+\tfrac{(1-2p)g''(p)}{2d}+O(d^{-2})$, so the large-$n$ bias is
\[
\mathrm{bias}=\gamma\big[\beta_x-(g(1)-g(0))\big]\;\xrightarrow{d\to\infty}\;\gamma\big[g'(p)-(g(1)-g(0))\big].
\]
\end{proposition}
\emph{Sketch.} In a $d$-regular graph without self-loops, $z_i$ and $x_i$ are uncorrelated with $\mathrm{Var}(z)=p(1-p)$, $\mathrm{Var}(x)=p(1-p)/d$, and bounded dependence gives convergence of $X^\top X/n$, so $\tilde c^\top(X^\top X/n)^{-1}\tilde c=(1,1)\,\mathrm{diag}(v,v/d)^{-1}(1,1)^\top=(1+d)/v$. For the slope, $x=B/d$ with $B\sim\mathrm{Bin}(d,p)$; expand $g$ around $p$ and use $E(x-p)^3=p(1-p)(1-2p)/d^2$. The direct effect $\tau_d$ is recovered since $z\perp x$. \hfill$\square$

\emph{Reading.} Under Bernoulli designs with many neighbours, nearly all units see $x_i\approx p$, so the regression identifies only the \emph{local} exposure--response slope at $p$; the GATE needs the secant slope from $0$ to $1$. They agree iff $g$ is (locally) linear, or by coincidence (quadratic $g$ at $p=\tfrac12$: tangent equals secant). Cluster randomisation spreads $x_i$ towards $0$ and $1$ (units in treated clusters see $x_i\approx1$), so the regression interpolates instead of extrapolating, at the price of correlation between $z$ and $x$ (larger $\|w\|$ if clusters are poor) and of fewer independent randomisation units. We test all of this empirically, including the proposition, in \S\ref{sec:res}.
